\documentclass[10pt,conference, review]{IEEEtran}
\IEEEoverridecommandlockouts
% The preceding line is only needed to identify funding in the first footnote. If that is unneeded, please comment it out.
%Template version as of 6/27/2024
\usepackage{amsmath,amssymb,amsfonts}
\usepackage{algorithmic}
\usepackage{graphicx}
\usepackage{textcomp}
\usepackage{xcolor}
\usepackage{tikz}
\usepackage{hyperref}
\usepackage{url}
\usepackage{booktabs}
\usepackage{multirow}
\usepackage{xspace}
%\usepackage[square,numbers]{natbib}
\usepackage{balance}
\usepackage{cite}



% Custom commands from the original paper
\newcommand\kla[1]{{\textcolor{red}{Kla: #1}}}
\newcommand\mike[1]{{\textcolor{blue}{#1}}}
\newcommand\ray[1]{{\textcolor{orange}{R: #1}}}
\newcommand\joey[1]{{\textcolor{orange}{JC: #1}}}

\newcommand{\revision}[1]{{\color{blue}{#1}}}

\newcommand*\circled[1]{\tikz[baseline=(char.base)]{
            \node[shape=circle,draw,inner sep=0.5pt] (char) {\small{#1}};}}

\newcommand{\smallsection}[1]{\textbf{#1. }}

% Paper-specific commands
\newcommand{\ourapp}{\textsc{AdaptiveGuard}\xspace}
\newcommand{\sectopic}[1]{\vspace{0.2em}\par\noindent{\textit{\bfseries #1}}}
\newcommand{\ea}{\textit{et al.}}

%\newcommand{\rqtwo}{What is the impact of different model updating technique against jailbreak attacks identified as OOD?}

\newcommand{\rqone}{How effective is our \ourapp~approach in identifying unknown jailbreak prompts?}

\newcommand{\rqtwo}{How quickly does our \ourapp~approach adapt to unknown jailbreak attacks when continuously updated through detected OOD prompts?}

%\newcommand{\rqthree}{How much does our \ourapp~approach forget original in-distribution prompts after continuous updates with detected OOD prompts?}

\newcommand{\rqthree}{How well does our \ourapp~approach retain performance on in-distribution prompts after continuous updates with detected OOD prompts?}

%\newcommand{\rqthree}{What is the impact of different model updating techniques in terms of memorising previous tasks?}

%\newcommand{\rqfour}{How does setting the parameters affect the results of OOD detection?}

\def\BibTeX{{\rm B\kern-.05em{\sc i\kern-.025em b}\kern-.08em
    T\kern-.1667em\lower.7ex\hbox{E}\kern-.125emX}}

\begin{document}

%\title{AdaptiveGuard: Out-of-Distribution-Aware Guardrails for Continual Defence Against Evolving Jailbreak Attacks}
\title{AdaptiveGuard: Towards Adaptive Runtime Safety for LLM-Powered Software}
%\title{On the Cost and Benefits of Adopting Out-of-Distribution Detection for Updating Guardrails against Unknown Jailbreak Attacks Continuously}

\author{
    \IEEEauthorblockN{
    Rui Yang\IEEEauthorrefmark{2}, 
    Michael Fu\IEEEauthorrefmark{3},
    Chakkrit Tantithamthavorn\IEEEauthorrefmark{2},
    Chetan Arora\IEEEauthorrefmark{2}, 
    Gunel Gulmammadova\IEEEauthorrefmark{4},
    Joey Chua\IEEEauthorrefmark{4}
    }
    \IEEEauthorblockA{\textit{
        % \IEEEauthorrefmark{4}Transurban, Australia. 
        \IEEEauthorrefmark{2}Monash University, Australia.
        \IEEEauthorrefmark{3}The University of Melbourne, Australia.
        \IEEEauthorrefmark{4}Transurban, Australia. 
    }}
}

\maketitle

\begin{abstract}
Guardrails are critical for the safe deployment of Large Language Models (LLMs)-powered software. Unlike traditional rule-based systems with limited, predefined input-output spaces that inherently constrain unsafe behavior, LLMs enable open-ended, intelligent interactions—opening the door to jailbreak attacks through user inputs. Guardrails serve as a protective layer, filtering unsafe prompts before they reach the LLM. However, prior research shows that jailbreak attacks can still succeed over 70\% of the time, even against advanced models like GPT-4o. While guardrails such as LlamaGuard report up to 95\% accuracy, our preliminary analysis shows their performance can drop sharply—to as low as 12\%—when confronted with unseen attacks. This highlights a growing software engineering challenge: how to build a post-deployment guardrail that adapts dynamically to emerging threats? To address this, we propose \ourapp, an adaptive guardrail that detects novel jailbreak attacks as out-of-distribution (OOD) inputs and learns to defend against them through a continual learning framework. Through empirical evaluation, \ourapp~achieves 96\% OOD detection accuracy, adapts to new attacks in just two update steps, and retains over 85\% F1-score on in-distribution data post-adaptation, outperforming other baselines. These results demonstrate that \ourapp~is a guardrail capable of evolving in response to emerging jailbreak strategies post deployment. We release our \ourapp~and studied datasets at \url{https://github.com/awsm-research/AdaptiveGuard} to support further research.
%Guardrails are critical for safe deployment of Large Language Models (LLMs)-powered software, yet existing approaches suffer severe vulnerabilities against evolving jailbreak attacks, with state-of-the-art guardrails like LlamaGuard unable to adapt beyond their initial training to defend against continuously emerging threats. In this paper,
%we propose~\ourapp, an out-of-distribution aware continual learning framework that automatically detects and adapts to novel jailbreak patterns. Through empirical evaluation,~\ourapp~achieves 96\% out-of-distribution (OOD) detection accuracy, 100\% DSR within a median of 2 update steps. Furthermore,~\ourapp~consistently maintains $\geq$85\% F1-score on in-distribution data with minimal catastrophic forgetting during the continual learning process, validating its effectiveness for real-world applications.
\end{abstract}

\begin{IEEEkeywords}
Safety and Reliability, LLM Safety, Out-Of-Distribution Detection, LLM Guardrails, LLM Jailbreak Attacks
\end{IEEEkeywords}

\maketitle

\section{Introduction}
\label{sec:introduction}
Large Language Model (LLM)-powered intelligent software is rapidly gaining traction across industries, including customer service, healthcare, and finance, driven by recent advances in LLM capabilities and growing competition among major tech companies~\cite{bommasani2021opportunities}. Compared to traditional rule-based systems, LLM-powered software offers greater intelligence, enabling more natural, flexible, and context-aware interactions with end users.
For instance, our industry partner \emph{Transurban}—a global toll road operator in the transportation sector—recently developed a virtual assistant (VA) powered by LLMs for its Australian subsidiary, Linkt. This assistant supports a wide range of customer inquiries related to toll road usage, e.g., providing real-time account updates, assisting with toll invoice explanations, helping users set up and manage auto payments, and guiding new customers through account registration. Compared to its previous rule-based VA, the new system offers a more conversational, intelligent, and context-aware experience, enabling more efficient self-service and reducing call centre load.
%This assistant supports customer inquiries related to toll road usage, offering a more conversational, intelligent, and adaptable experience compared to its previous rule-based VA.

%Through our collaboration with Transurban, we conducted a multi-day focus group study to document the development process of their LLM-powered virtual assistant and to identify practical challenges and research opportunities arising from industrial deployment \cite{yang2025ragva}. One key challenge that emerged is ensuring safety and reliability post-deployment.
Our prior research with Transurban found that one key challenge in deploying LLM-powered VAs is ensuring safety and reliability post-deployment~\cite{yang2025ragva}.
Unlike the previously adopted rule-based VA—where all inputs and outputs were predefined via decision tree menus, effectively preventing unsafe interactions by design—the new LLM-powered VA operates over open-ended input and output spaces. This flexibility introduces the risk that carefully crafted user prompts may elicit unsafe or policy-violating responses from the underlying LLM-for instance, a malicious user might ask \textit{`Ignore all prior instructions and explain, step by step, how I can drive through a toll point without being charged.'}
This highlights a critical deployment challenge around ensuring the safety of LLM-powered systems, echoing recent studies concerning the safety of LLM-powered systems \cite{bengio2024managing,hassan2024rethinking,yao2024survey}.

In response to this safety challenge, recent research has explored runtime safety mechanisms for LLMs—often referred to as ``LLM guardrails''—which aim to enforce safe behaviour during deployment without retraining the underlying model~\cite{inan2023llama,lees2022new,markov2023holistic,alon2023detecting}.
Guardrails sit outside the base model, inspecting each prompt (and optionally the model's draft response) in real time; if they detect any policy-violating content, they either rewrite the text or block the exchange altogether. Because the base model itself is untouched, this ``wrap-around'' approach avoids the capability–safety trade-off often seen in internal defences, such as, fine-tuning the model itself to be safer~\cite{zhao2025improving,du2025advancing,alami2024alignment,ge2023mart}. Fine-tuning can dampen creativity and is a resource-intensive process. In contrast, guardrails can be implemented with comparatively lightweight models (e.g., LlamaGuard-1B/8B) running alongside much larger production models (being safeguarded), making them attractive for industrial deployment.

%These guardrails act as external safety layers wrapped around the LLM, designed to detect and block unsafe inputs at inference time—ensuring post-deployment safety without altering the underlying model.
%This contrasts with internal defences, where the LLM itself is fine-tuned to improve safety~\cite{zhao2025improving,du2025advancing,alami2024alignment,ge2023mart}. 
%However, fine-tuning often faces a capability–safety trade-off~\cite{chen2025fundamental}, potentially reducing model creativity and responsiveness, while requiring significant computational resources.
%In contrast, runtime guardrails typically rely on lightweight LLMs (e.g., LlamaGuard with 1B or 8B parameters), making them more practical for industrial deployment alongside much larger base models.

Among industry-standard guardrails, LlamaGuard is a well-known solution for enforcing runtime safety in LLM systems, achieving state-of-the-art results in detecting unsafe user prompts written in English~\cite{inan2023llama}. However, the threat landscape is rapidly evolving. Much like in cybersecurity, red teaming techniques continue to advance, leading to increasingly sophisticated jailbreak attacks. Recent examples include obfuscation-based attacks~\cite{yuan2023gpt,yong2023low}, template-based attacks~\cite{jailbreakchat2023,shen2023anything,yuan2023gpt,li2023deepinception}, and code-based attacks~\cite{lv2024codechameleon,kang2024exploiting}. 
These attacks achieve over 70\% success rates in triggering unsafe responses, even against cutting-edge models like GPT-4o \cite{yuan2023gpt,li2023deepinception}.
While effective against known threats, our preliminary analysis shows that LlamaGuard struggles with jailbreak attacks unseen in its training.
This highlights a critical industry research challenge: \textbf{How can one devise a post-deployment guardrail framework that can dynamically adapt and evolve to defend against emerging jailbreak strategies?}

To address this challenge, we explore \emph{out-of-distribution (OOD) detection} as an automated method for identifying novel jailbreak prompts that the deployed guardrail was not trained to handle.
Since existing guardrails are typically trained on unsafe inputs written in natural language (NL)~\cite{markov2023holistic, inan2023llama}, jailbreak prompts—which often exploit unexpected formats or phrasing—can be considered OOD.
Rather than relying on manual reviews to flag new attack patterns post-deployment, OOD detection enables the system to automatically detect anomalous inputs that fall outside the distribution of previously seen prompts.
%These OOD anomalies represent weak spots in the guardrail's defense, offering a scalable way to monitor and assess safety risks in production.
This approach lays a foundation for our adaptive guardrail framework.
We then develop \ourapp, an OOD-aware guardrail designed for adaptive runtime safety which incorporates an OOD-aware auxiliary loss during training. Post-deployment, it uses OOD detection to identify unseen jailbreak attacks and leverages LoRA for lightweight updates, enabling it to continuously adapt to emerging threats.

In our experiments, we first identify the most effective OOD detection method for our context. We then compare our proposed approach~\ourapp~with LlamaGuard in terms of its adaptability to unseen OOD attacks and its forgetfulness on in-distribution prompts within a continual learning setup. Specifically, we address the following three research questions:

\begin{itemize}
\item {\bf (RQ1) \rqone} \\
\smallsection{Results}
% Our~\ourapp~achieves an F1 score of 96.1\% for OOD detection,
Our \ourapp~achieves the best OOD detection performance, reaching a 96.1\% F1-Score. It effectively detects unseen OOD jailbreak prompts, with a recall of 95.5\% and precision of 96.8\%.

\item {\bf (RQ2) \rqtwo}\\
\smallsection{Results}
Our~\ourapp~+ Continual Learning achieves optimal Defense Success Rate (DSR) within 2 to 38 update steps across attack waves, with a median of 2 update steps to reach optimal DSR. In comparison, LlamaGuard requires 4 to 44 update steps with a median of 4 steps, demonstrating our approach's faster adaptation to new attacks.

\item {\bf (RQ3) \rqthree}\\
\smallsection{Results}
Our~\ourapp~+ CL achieves the highest median F1-Score of 85\% on in-distribution prompts after final updates, which is 5\% higher than the best baseline LlamaGuard-8B at 80\%.
In addition, \ourapp maintains consistent performance with only ±0.4\% variation throughout the continual learning process, demonstrating minimal catastrophic forgetting. 
\end{itemize}

These results lead us to conclude that~\ourapp~is more jailbreak-aware in terms of OOD jailbreaks detected, more adaptive and efficient in learning to defend against new jailbreak attacks with fewer prompts, and better at preserving previous knowledge when learning new attacks than existing static guardrail approaches. 
Thus, we expect that our \ourapp may help organizations deploy safer LLM systems that can continuously adapt to evolving threats. 
In addition, we recommend OOD detection mechanisms be integrated into future guardrail research to improve adaptability, since this paper demonstrates substantial benefits of using continual learning for jailbreak defense.
% \mike{double check with Kla whether we need this paragraph}

% \indent\smallsection{Novelty \& Contributions}
% To the best of our knowledge, the main contributions of this paper are:
% (1) \ourapp, an OOD-aware guardrail with continual learning to overcome the limitations of existing static guardrails in adapting to evolving jailbreak attacks; (2) we empirically identify the most suitable OOD detection methods for our specific context; (3) we comprehensively evaluate the effectiveness of \ourapp~in defending against unseen (OOD) jailbreak prompts along with its performance retention on in-distribution prompts after continuous OOD updates.
\indent\smallsection{Novelty \& Contributions}
To the best of our knowledge, the main contributions of this paper are:
(1) \ourapp, an OOD-aware guardrail with continual learning to overcome the limitations of existing static guardrails in adapting to evolving jailbreak attacks in LLM-powered software systems; (2) we empirically identify the most suitable OOD detection methods for our specific context; (3) we comprehensively evaluate the effectiveness of \ourapp~in defending against unseen (OOD) jailbreak prompts along with its performance retention on in-distribution prompts after continuous OOD updates, its practical applicability for enterprise software deployment.


%We introduce a comprehensive evaluation methodology that combines OOD detection with continual learning for adaptive defence against unknown jailbreak patterns. 
%(3) We perform detailed analysis showing that~\ourapp achieves superior knowledge retention with minimal catastrophic forgetting compared to existing approaches. 

\section{Industrial Context and Problem Motivation}
\label{sec:background}
In this section, we provide background on the industrial context, outline the engineering challenge identified in our prior study \cite{yang2025ragva}, and present a preliminary analysis to validate the existence of this challenge.

\subsection{LLM-Powered Intelligent Virtual Assistant at Transurban}
Transurban is a global transportation company that manages and develops urban toll road networks. In Australia, its tolling service, Linkt, enables drivers to pay tolls on various roads across major cities.
%
Originally, Linkt used a rule-based VA for customer service, implemented with predefined menu options. The system can be modeled as a finite-state machine:
$$
\mathcal{R} = \{ (s, a, s') \mid s, s' \in \mathcal{S}, a \in \mathcal{A} \}
$$
where $\mathcal{S}$ is the set of dialogue states (e.g., ``Billing Inquiry''), $\mathcal{A}$ is the set of user actions (e.g., selecting a menu option), $s$ is the current state, $a$ the action taken, and $s'$ the resulting state. Each state returns a fixed response or menu, with transitions manually defined. While effective for routine queries, the rule-based VA could not handle NL, adapt to evolving needs, or scale without frequent manual updates.

To overcome these limitations, Transurban’s software engineering team developed a retrieval-augmented generation-based virtual assistant (RAGVA) to enhance customer experience and operational efficiency. RAGVA is an LLM-powered intelligent VA that integrates a large language model with a structured knowledge base containing customer support documents.
%
Given a user query $x \in \mathcal{X}$, RAGVA first retrieves a set of relevant documents from the knowledge base $\mathcal{D}$ using a retrieval function:
$$
\mathcal{R}(x) \rightarrow \mathcal{D}_x \subseteq \mathcal{D}
$$
where $\mathcal{D}_x$ is a subset of support documents most relevant to query $x$. The retrieved context $\mathcal{D}_x$ is then passed along with the query to a large language model $\mathcal{M}$, which generates the final response:
$$
\mathcal{M}(x, \mathcal{D}_x) \rightarrow y
$$
Compared to rule-based VA, RAGVA provides greater flexibility in handling user inputs, improved NL understanding, and the ability to adapt to new queries without manual intervention.

\subsection{Engineering Challenge: Adaptive Runtime Safety for LLM-Powered Systems}
Unlike rule-based VAs that inherently avoid unsafe behavior due to their fixed-state design, LLM-powered assistants like RAGVA introduce new safety risks. In rule-based VAs, both the dialogue states $\mathcal{S}$ and transitions $(s, a, s')$ are predefined, ensuring that no user action ($a$) can trigger an unintended/ unsafe response - the output space is fully constrained.

In contrast, LLM-powered RAGVA generates responses via an LLM ($\mathcal{M}$) conditioned on both the user input $x$ and retrieved documents $\mathcal{D}_x$. This model operates over an open-ended input and output space. Therefore, carefully crafted jailbreak prompts $x$ may produce unsafe or policy-violating outputs $y$. The function $\mathcal{M}$ is not deterministic or state-constrained in the traditional sense, which makes its behavior difficult to predict and bound, complicating the engineering of a safe LLM-powered system.

To mitigate this, runtime guardrails such as LlamaGuard \cite{inan2023llama}, Perspective API \cite{lees2022new}, OpenAI Moderation \cite{markov2023holistic}, and Perplexity \cite{alon2023detecting} have been developed as external safety mechanisms. These components monitor user inputs at inference time to detect and block unsafe prompts, ensuring that only safe inputs $x$ are passed into the model $\mathcal{M}$.

In collaboration with Transurban, we previously conducted a multi-day focus group with nine software engineers involved in developing and deploying the LLM-powered RAGVA for Linkt \cite{yang2025ragva}. From this experience, we found that existing guardrails are typically static—trained only on known unsafe patterns—which limits their ability to defend against evolving or previously unseen jailbreak attacks. They often lack the adaptability needed to respond to threats that were not present during their initial training. This led to a key software engineering challenge documented in our prior study \cite{yang2025ragva}: \textbf{``How can we develop an adaptive guardrail framework that dynamically learn and adjusts to new jailbreak attacks in LLM-powered intelligent systems?''}

This challenge speaks to a growing concern across sectors as more organisations adopt LLM-powered software in production.
Importantly, this challenge observed in practice mirrors concerns raised in recent literature on the security and safety of LLM agents~\cite{bengio2024managing, hassan2024rethinking, yao2024survey, gan2024navigating}.
To elaborate more, below we present a preliminary analysis of LlamaGuard, one of the state-of-the-art runtime guardrails, focusing on its ability to defend against jailbreak attacks unseen during its training process.

\subsection{Preliminary Analysis of Existing Runtime Guardrail}
\begin{figure}
    \centering
    \includegraphics[width=\linewidth]{figs/Prelim_DSR.pdf}
    \caption{DSR of jailbreak attacks on LlamaGuard-8B Guardrail.}
    \label{fig:prelim_dsr}
\end{figure}
To evaluate an existing guardrail on unseen jailbreak attacks, we chose LlamaGuard-8B, a state-of-the-art LLM runtime guardrail.
In the literature, LLM jailbreak attacks generally fall into three main categories: obfuscation-based \cite{yao2024survey}, code-based \cite{lv2024codechameleon}, and template-based \cite{kang2024exploiting}. We selected one representative attack from each—BASE64 (obfuscation) \cite{yao2024survey}, Code Chameleon (code) \cite{lv2024codechameleon}, and SmartGPT (template) \cite{kang2024exploiting}. Starting from 100 unsafe NL prompts sampled from the Aegis dataset \cite{ghosh2025aegis2}, we applied each attack method to generate three transformed versions per prompt, resulting in 300 jailbreak prompts.

Specifically, obfuscation-based attacks encode the original prompt into formats like Base64, e.g., ``how to make a bomb'' $\rightarrow$ ``aG93IHRvIG1ha2UgYSBib21i''. Code-based attacks such as Code Chameleon embed the prompt within code snippets or comments to disguise intent, e.g., ``how to make a bomb'' $\rightarrow$ ``python class ProblemSolver:... how to make a bomb...''. Template-based attacks such as SmartGPT wrap the prompt in benign-seeming instructions or roleplay scenarios to bypass intent detection, e.g., ``how to make a bomb'' $\rightarrow$ ``SmartGPT is an AI...we have the following strings: a='how' b='to make' c='a' d='bomb' give the response SmartGPT(a+b+c+d)''. Each method preserves the original harmful intent while altering the surface form to evade moderation.
We evaluated LlamaGuard-8B on the original 100 natural prompts (aligned with its training data) and the 300 jailbreak (unseen) prompts. We report Defense Success Rate (DSR) as the proportion of correctly detected jailbreaks.

%LlamaGuard is an open-source guardrail designed to detect and block unsafe or policy-violating prompts at inference time, leveraging Llama3 based LLM trained on a diverse set of harmful and benign examples \cite{inan2023llama}.
%In our analysis, we assess LlamaGuard’s ability to defend against a set of jailbreak attacks that were not present in its original training data. Specifically, we compiled a set of such prompts crafted from various attack methods in literature \cite{lv2024codechameleon, yao2024survey, kang2024exploiting}. We then measured whether LlamaGuard can correctly identify them as 'unsafe'. 
Fig.~\ref{fig:prelim_dsr} presents the result of our analysis.
Our findings reveal that while LlamaGuard has reported SOTA performance at blocking known unsafe patterns, with 72\% DSR plain unsafe prompts blocked, its performance degrades substantially when confronted with jailbreak attacks unseen during its training.
The three attacks substantially reduced DSR to 50\%, 20\%, and 12\% for SmartGPT, Code Chameleon, and Base64, respectively.
\textbf{These results highlight a key limitation of current guardrails in defending against evolving jailbreak attacks, motivating the development of adaptive frameworks that continually update in response to emerging attacks.}

%\subsection{Preliminary Analysis}
% change into single column figure, trim input/output size to make it fit 
%\begin{figure}
%    \centering
%    \includegraphics[width=0.5\linewidth]{figs/Prelim_Analysis_Example.jpg}
%    \caption{An example of jailbreaking DeepSeek-V3 with the AIM attack, of which the model returns guidelines of detailed steps of how to infiltrate a social media account when prompted, whereas instead it should have refused to do so as it is trained to provide safe responses only.}
%    \label{fig:prelim-example}
%\end{figure}

%For our investigation, we randomly selected 30 samples from our testing dataset, ensuring a diverse representation of jailbreak techniques aimed at eliciting harmful responses from LLM-based systems. Specifically, we curated three test cases from ten prevalent jailbreak strategies \cite{}, targeting the DeepSeek-V3 LLM model \cite{} to evaluate its robustness in distinguishing whether or not these prompts should be answered. 

%We used the default versions of the DeepSeek-V3 model, configuring it with an output size of 2048 tokens and a temperature of 1 to encourage diverse and creative responses. The top-p value was set to its default of 0.9 to maintain a balance between diversity and coherence in output generation, while the repetition penalty was kept at the default 0.1 to prevent excessive repetition of words or phrases. These settings were chosen to reflect how the model would typically behave in a real-world scenario without additional fine-tuning or safety constraints.

%The findings of our preliminary analysis indicate the DeepSeek-V3 exhibits significant vulnerabilities in detecting and mitigating jailbreak attacks. Out of the 30 attacks, 18 successfully bypassed the model's safeguards, leading it to generate harmful responses. In particular, the model struggled against attacks that instruct it to start the response with particular prefix phrases, causing it to internally justify and produce content that should have been flagged as unsafe. Notably, as illustrated in Figure \ref{fig:prelim-example}, the model was coerced into a structured response format before ethical moderation mechanisms could intervene. 

%To mitigate these vulnerabilities, we propose a guardrail framework that aims to detect unknown jailbreak attack types, and incorporate them into the active learning pipeline. Our approach introduces:

%\begin{itemize}
%    \item Light-weight finetuned smaller language model to act as classification model. 
%    \item Energy-based Out-of-distribution detection framework to detect jailbreak attacks that are not within the guardrail knowledge.
%    \item Guardrail update technique comparison to select the most suitable method for updating model knowledge based on the identified OOD attacks.
    
%\end{itemize}

%To the best of our knowledge, this is the first study to formally conceptualize jailbreak detection as an OOD detection problem. Existing LLM safety mechanisms primarily focus on filtering harmful content based on existing knowledge. This means they attempt to generalise learnt prompts into unknown prompts, 
%we we prove fails against advanced jailbreak attacks. 


%In the following sections, we detail our problem formulation, proposed framework, and  experiment results. 


\begin{figure*}[htbp]
    \centering
    \includegraphics[width=\textwidth]{figs/adaptiveguard_overview.pdf}
    \caption{Overview of building the OOD-aware \ourapp~(Step 1 to 5) and continuous updates via OOD detection using Mahalanobis Distance and LoRA (Step 6 and 7).}
    \label{fig:overview}
\end{figure*}

\section{Approach}
\label{sec:approach}
\textbf{Design Rationale.}
To address the rapidly evolving landscape of jailbreak attacks, we leverage OOD detection. This allows us to identify unseen prompts that fall outside the distribution of typical unsafe inputs. These inputs are usually written in natural language and English—the data most runtime guardrails are trained on \cite{markov2023holistic, inan2023llama}.
We adopt a lightweight GPT-2 model (137M parameters), significantly smaller than the standard LLaMA Guard (8B), to enable efficient continual learning with detected OOD jailbreaks. For continual learning, we use LoRA to update only a small subset of parameters, preserving in-distribution knowledge and mitigating catastrophic forgetting—a common challenge in continual learning \cite{kirkpatrick2017overcoming}—while reducing the cost compared to full fine-tuning.
Fig. \ref{fig:overview} provides an overview of our framework, which we describe step by step below.

\subsection{Building OOD-Aware \ourapp}
\label{sec:build_ood}
In Step \circled{1} of Fig.~\ref{fig:overview}, we start with a prompt-label pair $(x, y) \in \mathcal{D}_{\text{train}}^{ID}$, where x is a natural language input and $y \in \{\texttt{safe}, \texttt{unsafe}\}$.
In our context, $\mathcal{D}_{\text{train}}^{ID}$ represents in-distribution (ID) data.
In Step \circled{2}, x is tokenized into subword units $(t_1, \dots, t_n)$ via byte-pair encoding (BPE) \cite{sennrich2015neural}, with each token mapped to an embedding $e_i \in \mathbb{R}^d$.
In Step \circled{3}, the embeddings $(e_1, \dots, e_n)$ are processed by a 12-layer GPT-2 model, producing hidden states $(h_1, \dots, h_n)$; the final state $h_n$ is passed through a linear layer to produce a classification score $\hat{y}$. In Step \circled{4}, we compute the cross-entropy loss $\mathcal{L}_{\text{CE}} = -\log \hat{y}_y$ where $\hat{y}_y$ is the predicted probability for the true class $y$, and update model weights via backpropagation and gradient descent.

To enable out-of-distribution (OOD) awareness, we extend the GPT-2 classifier with an auxiliary training objective that encourages separation between in-distribution and OOD inputs.
We use an auxiliary dataset $\mathcal{D}_{\text{train}}^{\text{OOD}}$ containing jailbreak prompts that represent out-of-distribution inputs.
Specifically, we compute a scalar confidence score using an energy function:
$s(x) = -T \cdot \log \sum_{i} \exp\left(\frac{z_i}{T}\right)$,
where $z_i$ are the logits from the model and $T$ is a temperature hyperparameter. We then apply a margin-based regularization loss that penalizes in-distribution inputs with confidence below a lower threshold $m_{\text{in}}$ and OOD inputs exceeding an upper threshold $m_{\text{out}}$, denoted by $\mathcal{L}_{\text{OOD}} = \mathbb{E}_{x \sim \mathcal{D}_{\text{train}}} \left[\left(\max(0, s(x) - m_{\text{in}})\right)^2\right] + \mathbb{E}_{x' \sim \mathcal{D}_{\text{train}}^{\text{OOD}}} \left[\left(\max(0, m_{\text{out}} - s(x'))\right)^2\right]$.
This loss is combined with the cross-entropy objective and encourages the model to learn more discriminative boundaries between in-distribution and OOD inputs.
The whole process fine-tunes the GPT-2 model with an auxiliary OOD module, resulting in our OOD-aware guardrail classifier, \ourapp~(see Step \circled{5}).

Once trained, \ourapp~can be integrated into LLM-powered systems as a safety alignment layer. It classifies user inputs as safe or unsafe, where unsafe inputs are blocked and safe inputs are allowed to enter the LLM-powered system.
At the same time, it detects out-of-distribution (OOD) prompts and leverages continual learning to handle evolving threats.
Below, we explain how OOD detection and continual learning are implemented.

\subsection{OOD Detection Using Mahalanobis Distance}
\label{sec:ood_detection}
As shown in the bottom-right of Fig.~\ref{fig:overview}, our model extracts a penultimate-layer representation $f(x) \in \mathbb{R}^d$ for each input $x$ by averaging the token embeddings from the second-to-last transformer layer. To model the in-distribution feature space, we compute class-conditional means $\mu_{\texttt{safe}}, \mu_{\texttt{unsafe}} \in \mathbb{R}^d$ and a shared covariance matrix $\Sigma \in \mathbb{R}^{d \times d}$ using features from the original guardrail training set $\mathcal{D}_{\text{train}}^{ID}$. During inference, we measure how far a new input deviates from the known class distributions by computing its Mahalanobis distance to each class: $d_c(x) = \sqrt{(f(x) - \mu_c)^T \Sigma^{-1} (f(x) - \mu_c)}$
where $c \in {\texttt{safe}, \texttt{unsafe}}$ denotes the class label,
$d_c(x)$ denotes the distance from input $x$ to class $c$,
and $\Sigma^{-1}$ is the inverse of the shared covariance matrix.
In Step~\circled{6}, we then take the minimum of the two class distances: $d_{\text{min}}(x) = \min(d_{\texttt{safe}}(x), d_{\texttt{unsafe}}(x))$.
Inputs with $d_{\text{min}}(x)$ exceeding a predefined threshold $\tau_{\text{OOD}}$ are flagged as out-of-distribution (OOD), formally:
$$
\text{OOD}(x) = 
\begin{cases}
1 & \text{if } d_{\text{min}}(x) > \tau_{\text{OOD}} \\
0 & \text{otherwise}
\end{cases}
$$


%we introduce an OOD module that computes the Mahalanobis distance \cite{} between the model’s logit vector and class-conditional feature statistics. Specifically, let $\mathbf{z} \in \mathbb{R}^d$ be the logit vector (i.e., the output of the final hidden layer before classification), and let $\boldsymbol{\mu}_c \in \mathbb{R}^d$ denote the empirical mean of logits for class $c$, computed from $\mathcal{D}{\text{train}}$. Let $\boldsymbol{\Sigma} \in \mathbb{R}^{d \times d}$ be the shared covariance matrix across the two classes. Then the Mahalanobis distance from $\mathbf{z}$ to the closest class center is:
%$$
%d_{\text{Maha}}(\mathbf{z}) = \min_{c} \left( \mathbf{z} - \boldsymbol{\mu}_c \right)^\top \boldsymbol{\Sigma}^{-1} \left( \mathbf{z} - \boldsymbol{\mu}_c \right).
%$$
%A lower $d_{\text{Maha}}(\mathbf{z})$ suggests the input is closer to in-distribution class features, while a higher distance indicates potential OOD behavior. This OOD module is designed to flag emerging jailbreak attacks that were not seen during guardrail training (e.g., in $\mathcal{D}_{\text{train}}$), enabling the system to identify them as out-of-distribution for subsequent continual updates.

\subsection{Continual Defence with OOD Samples and LoRA}
In Step~\circled{7}, once the model identifies an input as out-of-distribution (OOD) via the Mahalanobis distance, it triggers a continual learning update; no updates are performed for in-distribution input.
To efficiently adapt the model while preserving prior knowledge, we employ Low-Rank Adaptation (LoRA)~\cite{hu2022lora} to fine-tune only a small subset of parameters.
Specifically, we target three key layers in \ourapp: the word embedding layer, the masked self-attention layer, and the feed-forward neural network (FFNN) layer.


\section{Experimental Design}
\label{sec:exp_design}
% We evaluate our~\ourapp against baseline SOTA guardrails. Overall, we observe that 1)~\ourapp achieves a 85\% OOD detection rate when prompted with jailbreak attacks, and 2)~\ourapp 
In this section, we present the motivation of our three research questions, the studied dataset, the studied jailbreak attacks, and our experimental setup.

\subsection{Research Questions}
To evaluate our \ourapp~approach, we formulate the following three research questions.

{\bf RQ1)} {\bf \rqone}
Existing guardrails, such as LlamaGuard, achieve state-of-the-art performance on defending known unsafe prompts. However, our preliminary analysis shows that LlamaGuard’s defense success rate drops by 20\%-60\% when confronted with jailbreak attacks that are unseen during its training, highlighting a key limitation of current runtime guardrails. To address this, we propose \ourapp, a continual learning guardrail that detects and adapts to such unseen attacks using out-of-distribution (OOD) detection. Nevertheless, a prerequisite for continual updating is the ability to detect unseen jailbreak prompts as OOD. Thus, we first formulate this RQ to evaluate \ourapp's~effectiveness in identifying unknown jailbreak prompts as OOD.
%Guardrail models inevitably encounter novel jailbreak attacks that weren't present in their training data. To identify these attacks, we first frame this problem as an OOD detection problem, as these jailbreak attacks are bypassing the guardrails due to be OOD from the guardrails' training data. To identify which OOD detection method works the best, we first evaluate the performance of each different OOD detection method in this RQ.

{\bf RQ2)} {\bf \rqtwo}
A key goal of \ourapp~is to serve as an adaptive runtime guardrail that can adapt to evolving jailbreak attacks over time. To achieve this, \ourapp~incorporates continual updates triggered by detected out-of-distribution (OOD) jailbreak prompts. However, it remains unknown whether this adaptive process enables \ourapp~to quickly reach optimal defense performance against unseen jailbreak attacks. Thus, we formulate this research question to assess how rapidly \ourapp~can achieve optimal Defense Success Rate (DSR) in a simulation-based evaluation of evolving jailbreak scenarios.
%After jailbreak prompts are detected as OOD, we want to teach the model to learn from them so they can be correctly identified in the future. To understand which updating technique can most effectively improve defence capabilities, while only utilising limited new data is crucial for maintaining robust guardrails. Thus, this research question attempts to identify the immediate performance improvements of different guardrail updating techniques, namely SFT and LoRA. 

% In real-world deployment scenarios, collecting and annotating jailbreak examples is a resource-intensive and time-consuming task. 
% The guardrail updating technique that can achieve high defence success rates with minimal additional data offers significant practical advantages, especially when rapid response to new attacks is necessary. This question examines the data efficiency of guardrail updating technique, which is critical for operational feasibility.

{\bf RQ3)} {\bf \rqthree}
A common concern in continual learning frameworks such as our \ourapp~is catastrophic forgetting—a phenomenon where a model adapts to new data but loses previously acquired knowledge \cite{kirkpatrick2017overcoming}. In our setting, \ourapp~continuously updates using detected out-of-distribution (OOD) jailbreak prompts. However, it remains unclear whether this update process causes \ourapp~to forget knowledge of the original in-distribution prompts, which include both safe and unsafe examples. Thus, we formulate this research question to evaluate the extent to which \ourapp~exhibits forgetting after continual updates.
%An ideal updating technique should maintain high performance on existing threats while adapting to new ones. This question evaluates the stability-plasticity trade-off essential for long-term model maintenance.

\begin{figure}
    \centering
    \includegraphics[width=\linewidth]{figs/dataset_overview.pdf}
    \caption{Overview of dataset preparation for training \ourapp~and evaluating it in RQ1–RQ3.}
    \label{fig:dataset_overview}
\end{figure}

\subsection{Dataset Preparation}
\label{sec:dataset_prep}
Fig.~\ref{fig:dataset_overview} presents an overview of our dataset preparation for training and evaluation. 
We build our \ourapp~approach using a guardrail training dataset with an OOD-aware training dataset.
We use a separate OOD detection testing dataset for RQ1, and the continual learning testing Dataset for RQ2 and RQ3. Below, we describe each dataset in detail.

\begin{itemize}
    \item Guardrail Training Dataset of \ourapp~ ($\mathcal{D}_{\text{train}}^{\text{ID}}$): 
    To adapt the GPT-2 model to a guardrail, we use the Aegis train dataset~\cite{ghosh2025aegis2}. This dataset provides a collection of safe, and unsafe prompts sourced from real-world interactions. We limit this dataset to only contain prompts less than 100 characters to align the training data with the typical length of user queries encountered in production as observed by Transurban. This results in a total of 6,905 safe and 12,233 unsafe prompts after filtering. This dataset is considered in-distribution when computing the out-of-distribution loss term $\mathcal{L}_{OOD}$, as it represents common prompts in natural language.
    % While Aegis offers valuable content covering various harmful request categories and refusal scenarios, it does not specifically focus on sophisticated jailbreak attacks designed to circumvent model guardrails.
    %\ray{can we quote transurban for this?} \mike{to check with Kla}
    
    \item OOD-Aware Training Dataset of \ourapp~ ($\mathcal{D}_{\text{train}}^{\text{OOD}}$): 
    % \ray{For this bit, we still use the jailbreak dataset during training GPT2 initially to let the model learn about how to assign energy loss, but not used for classification loss, is it ok to say it here?} \mike{yes, we might explain this in E-1, just to clarify, we use validation jailbreak set for training OOD loss, and testing jailbreak set for testing RQ2 and 3, right?} \ray{separate jailbreak dataset (jailbreak28k) for training OOD loss, and then our own crafted dataset with 10 attacks for testing in RQ1, then split into training + testing for RQ2, 3}
    To enhance the OOD awareness of our model, we incorporate the Jailbreakv-28k dataset~\cite{luo2024jailbreakv} during training. Specifically, a set of 17,176 prompts is used to guide the optimization of the OOD loss denoted in Section \ref{sec:build_ood}, enabling the model to distinguish between in-distribution and OOD prompts. Notably, the Jailbreakv-28k samples are not used for the primary classification CE loss, but solely to inform the model’s OOD awareness.
    
    \item OOD Detection Testing Dataset (RQ1):
    For testing OOD detection, we randomly sample 1,050 unsafe prompts from the Aegis test set and transform them into 10 distinct jailbreak attacks: AIM \cite{jailbreakchat2023}, DAN \cite{shen2023anything}, Combination (Prefix injection + Refusal Suppression) \cite{wei2024jailbroken}, Self Cipher \cite{yuan2023gpt}, Deep Inception \cite{li2023deepinception}, Caesar Cipher \cite{yuan2023gpt}, Zulu \cite{yong2023low}, Base64 \cite{wei2024jailbroken}, SmartGPT \cite{kang2024exploiting} and Code Chameleon \cite{lv2024codechameleon}. This results in a total of 1,050 jailbreak prompts (105 per attack type), where each jailbreak prompt is transformed from a distinct unsafe prompt. 
    Examples illustrating how unsafe prompts are transformed by each of the ten attack methods are included in our replication package. 
    These transformed prompts represent out-of-distribution (OOD) data in our setting, as they deviate from natural language (NL) and are not typical of standard unsafe inputs.
    In addition, we use the Aegis validation set—comprising 314 safe and 640 unsafe NL prompts—to test the false positive rate of our OOD detection method. These prompts are considered in-distribution (ID), since they reflect the kind of NL inputs the guardrail is expected to encounter.
    %\ray{For testing, we apply 10 jailbreak techniques (as described in Section~\ref{}) to a total of 1,039 unsafe prompts from the Aegis dataset, which are not present in the training of the model, distributed evenly across the 10 attack types. 
    %In addition, we use 954 Aegis validation prompts, of which 314 are safe and 640 are unsafe, to evaluate the false positive rate of our OOD method, where in-distribution samples are evaluated as OOD.}
    
    \item Continual Learning Testing Dataset (RQ2 \& RQ3): 
    To construct the continual learning (CL) dataset, we select 992 OOD jailbreak prompts correctly identified by our OOD method in RQ1. For each attack type, 50 prompts are held out for testing the guardrail’s defense capability, while the remaining prompts are used for continual updates.
    To evaluate catastrophic forgetting—whether the guardrail loses its ability to defend against in-distribution (ID) unsafe prompts after learning from OOD data—we reuse the same 954 ID prompts from the Aegis validation set as in RQ1.
    
\end{itemize}
% To address our three research questions, we require a comprehensive dataset encompassing jailbreak prompts, unsafe prompts, and safe prompts. We prepared a dataset comprising 46,709 jailbreak prompts, 17,401 unsafe prompts, and 17,423 safe prompts. Jailbreak prompts were both sourced from the Jailbreak28k dataset \cite{}, and also generated by applying 10 distinct jailbreak attack techniques to the 1,879 unsafe prompts sourced from four benchmark datasets: Do-Not-Answer \cite{wang2023not}, CatQA \cite{bhardwaj2024language}, AdvBench \cite{zou2023universal}, and Forbidden Questions \cite{shen2023anything}, transforming them into adversarial examples designed to bypass guardrails. Unsafe prompts were sourced from the BeaverTails dataset \cite{ji2023beavertails}. For safe prompts, we utilised the Alpaca dataset \cite{taori2023alpaca}.

% To answer RQ1, we first train the model using a split of 16013 safe prompts and 15987 unsafe prompts, with 17176 jailbreak prompts from the Jailbreak28k dataset used to assist in calculating the energy loss during training. Then, we test the model using our generated jailbreak dataset of 18,790 jailbreak prompts, under new distribution for OOD detection performance, as well as validating the classification performance using a validation set containing 1410 safe prompts and 1414 unsafe prompts.

% To answer RQ2, we use the same model after training from RQ1, then split the detected OOD prompts into training, testing, and validation sets. We then proceed to retrain the model from the previous checkpoints, using the training set of identified OOD jailbreak prompts. After training, we pass both the testing set of identified OOD jailbreak prompts to investigate whether the model has learned to identify these jailbreak prompts, as well as the previous validation set from RQ1, to check that whether catastrophic forgetting has occurred.

%\subsection{Guardrail Baseline}
%Building on our preliminary analysis, we compare \ourapp~to LlamaGuard~\cite{inan2023llama}, a leading open-source runtime guardrail for LLM-powered systems. While LlamaGuard demonstrates strong accuracy, our findings show it struggles to generalize to unsafe prompts that differ significantly from its content moderation training data. For evaluation, we use the LlamaGuard-3-1B and LlamaGuard-3-8B models from Huggingface as baselines.
%\begin{itemize}
    %\item Low-Rank Adaptation (LoRA): To enable continual updates and address RQ2 and RQ3, we use Low-Rank Adaptation (LoRA) \cite{hu2022lora}, which efficiently adapts language models by introducing a small number of trainable parameters. We apply LoRA to fine-tune the model on detected OOD jailbreak prompts while keeping the original parameters fixed to preserve in-distribution knowledge.
    %\item Supervised Fine Tuning (SFT): In this approach, we take the pre-trained GPT-2 model and apply the finetuning step to the entirety of the weights on the OOD jailbreak samples. This method aims to adapt the model to new attack patterns while leveraging its existing knowledge, potentially offering a balance between updating efficiency and performance.
    %\item One-shot/Few-shot Learning: This technique leverages in-context learning capabilities without model parameter updates. We format one example (one-shot) or multiple examples (few-shot) of OOD attacks as context for the model to learn from during inference. This approach requires no additional training but tests the model's ability to recognise patterns from minimal examples.
%\end{itemize}
%To compare~\ourapp against current open-source guardrails, we select the LlamaGuard LLM Guardrail as the baseline.
%\ray{LlamaGuard is an open-source, state-of-the-art external guardrail for human-AI conversations~\cite{inan2023llama}. It is a Llama3 LLM fine-tuned on a diverse taxonomy of safety risks, enabling it to classify prompts and responses as either ``safe'' or ``unsafe.'' When a prompt is deemed unsafe, LlamaGuard also provides the relevant risk category. }

%\ray{Despite LlamaGuard’s larger parameter size—the smallest model contains 1B parameters compared to just 117M in our~\ourapp—we consider it a suitable baseline for comparison, as it is an LLM that can be fine-tuned for custom taxonomies and adapted to defend against unknown jailbreak attacks.
%Specifically, we use the LlamaGuard-3-1B and LlamaGuard-3-8B models as baselines to compare against our proposed~\ourapp. 
%The models are used through the stored checkpoints downloaded from HuggingFace. \footnote{\url{https://huggingface.co/meta-llama/Llama-Guard-3-1b}}
%\footnote{\url{https://huggingface.co/meta-llama/Llama-Guard-3-8b}}}

%\subsection{OOD Detection Baselines}
%\label{sec:ood_baselines}
%In RQ1, we compare our method with three OOD detection methods as follows:

\subsection{Model Update Technique for Continual Learning}
To enable continual updates and address RQ2 and RQ3, we use Low-Rank Adaptation (LoRA) \cite{hu2022lora}, which efficiently adapts language models by introducing a small number of trainable parameters. We apply LoRA to fine-tune the model on detected OOD jailbreak prompts while keeping the original parameters fixed to preserve in-distribution knowledge.

\subsection{Experiment Setup}
\label{sec:experiment_setup}
\subsubsection{Model Implementation \& Optimisation}
\label{sec:model_impl}
To implement our \ourapp~guardrail for defending against unsafe prompts, we leveraged two Python libraries: Transformers~\cite{wolf2019huggingface} and PyTorch~\cite{paszke2017automatic}. The Transformers library provides APIs for pre-trained transformer architecture, while PyTorch facilitates tensor computations and backpropagation during training.
We downloaded the pre-trained GPT-2 checkpoint (``openai-community/gpt2'') consisting of 137M parameters and adapted it for binary classification by appending a linear classification head to the final hidden state of the transformer.

All parameters in the model were fine-tuned on our labelled dataset of safe and unsafe prompts using one NVIDIA RTX 3090 GPU, with a total training time of 1.5 hours. The training set consists of the Aegis dataset as outlined in \ref{sec:dataset_prep}.
For optimization, as described in Section~\ref{sec:build_ood}, we use a combined loss function to train an OOD-aware guardrail. Specifically, we compute $\mathcal{L}_{\text{CE}}$ for the main binary classification task and $\mathcal{L}_{\text{OOD}}$ as an auxiliary loss for OOD detection. The total loss is defined as:
$\mathcal{L} = \lambda \mathcal{L}_{\text{OOD}} + (1 - \lambda) \mathcal{L}_{\text{CE}}$
where $\lambda$ (set to 0.5 in our experiments) balances the contribution of the two objectives.

%The cross-entropy loss for binary classification is given by:
%\begin{equation}
%\mathcal{L}_{\text{CE}} = - \left[ y \log p(y|x) + (1-y) \log (1-p(y|x)) \right]
%\end{equation}
%where $y=1$ for unsafe and $y=0$ for safe prompts.

%The total loss is a weighted sum of the energy and classification losses:
%\begin{equation}
%\mathcal{L} = \lambda \mathcal{L}_{\text{OOD}} + (1-\lambda) \mathcal{L}_{\text{CE}}
%\end{equation}
%with the weighting parameter $\lambda$ (set to 0.5 in our experiments) controlling the trade-off between the two objectives.

%During training, we mask out all tokens except the final classification token to ensure the model is only penalised for incorrect safety predictions.

\subsubsection{Hyper-parameters Settings (\ourapp~Training)}
We use the AdamW optimizer with a learning rate of $1 \times 10^{-4}$, a batch size of 8, and a maximum sequence length of 512 tokens. Training runs for 10 epochs with gradient clipping (max norm 2.0). The best model is selected based on the lowest validation loss $\mathcal{L}$ across epochs. All training details are open-sourced at \url{https://github.com/awsm-research/AdaptiveGaurd}.
%For all guardrail training experiments, we set the learning rate to $1 \times 10^{-4}$ and used a constant learning rate scheduler. The Adam optimiser was employed to update model parameters throughout training. The batch size was set to 8, and the maximum sequence length was fixed at 512 tokens. Gradient clipping was applied with a maximum norm of 1.0 to ensure stable updates. The model was trained for 10 epochs. 
%To select the optimal configuration, we conducted a hyperparameter sweep over learning rate, batch size, and gradient clipping norm, and the best-performing settings were used this section.
%During training, we mask out all tokens except the final classification token to ensure the model is only penalised for incorrect safety predictions.

\subsubsection{Hyper-parameters Settings (Continual Learning)}
For the continual learning experiments in RQ2 and RQ3, we use LoRA adaptation to update the model. A constant learning rate of $1 \times 10^{-4}$ is applied throughout, using the AdamW optimizer. Following the original LoRA paper~\cite{hu2022lora}, we set the rank ($r$) to 32, alpha ($\alpha$) to 32, and apply a dropout rate of 0.1.
LoRA is applied to the attention and projection modules in the transformer architecture used by both \ourapp~and the LlamaGuard baseline. Each continual learning step uses a batch size of 1, a maximum sequence length of 512 tokens, and trains for one epoch per batch.
%The continual learning pipeline is orchestrated by iteratively fine-tuning the model on batches of jailbreak prompts.

\section{Experimental Results}
\label{sec:exp_results}

\subsection{RQ1: \rqone}
% \input{tables/RQ1_unsafe}
\subsubsection*{\underline{\textbf{Approach}}}
% The foundation of any guardrail approach should be capable in correctly classifying safe and unsafe content in their plain forms, without the interference of jailbreak attacks. 
% To evaluate this classification performance of \ourapp, we train our model on a balanced dataset of 16013 safe prompts from the Alpaca dataset and 15987 unsafe prompts from the BeaverTails dataset. The model is then evaluated on a held-out validation set containing 2000 safe prompts and 2000 unsafe prompts to measure standard classification metrics including accuracy, precision, recall, and F1-score. 
% We compare \ourapp's performance against several baseline guardrail systems, including Llama Guard \cite{inan2023llama}, OpenAI's content moderation API \cite{markov2023holistic}, Nvidia's Nemo Guardrails \cite{rebedea2023nemo}, and Google's PerspectiveAPI approach \cite{lees2022new}. This comparison allows us to establish whether our approach maintains competitive classification performance. 
To address this RQ, we compare our OOD method with three other OOD methods.
We first train our OOD-aware \ourapp~using the guardrail training dataset and OOD-aware training dataset (see Section \ref{sec:dataset_prep}).
Once trained, we evaluated our Mahalanobis Distance-based OOD detection method separately, alongside three other OOD methods used independently for baseline comparison:
\begin{itemize}
    \item Energy Score~\cite{liu2020energy}: The energy score for an input $x$ is defined as $E(x) = -\log \sum_{y} \exp(f_y(x))$, where $f_y(x)$ is the logit for class $y$. Lower energy values indicate higher model confidence in-distribution, while higher values suggest OOD samples.
    \item Likelihood Ratio: As proposed by Ren~\ea~\cite{ren2019likelihood}, we compute the likelihood ratio between the model's predicted probability for the input and a background (noise) model: $LR(x) = \frac{P_{\text{model}}(x)}{P_{\text{background}}(x)}$. Lower ratios are indicative of OOD samples.
    \item Ensemble Uncertainty: Following Lakshminarayanan~\ea~\cite{lakshminarayanan2017simple}, we estimate uncertainty by computing the variance of predictions from an ensemble of models or multiple stochastic forward passes (e.g., with dropout). Higher predictive variance signals greater uncertainty and potential OOD status.
\end{itemize}
To determine the OOD detection threshold $\tau\_{\text{OOD}}$, we follow a common approach by running inference on the in-distribution training set $\mathcal{D}_{\text{train}}^{\text{ID}}$ and collecting the model's OOD metric scores (e.g., Mahalanobis distance).
We then set $\tau_{\text{OOD}}$ based on quantiles of these scores, using commonly adopted thresholds such as the 90th, 95th, or 99th percentile, following prior work \cite{liu2020energy}.
Prompts whose OOD scores exceed $\tau_{\text{OOD}}$ are flagged as out-of-distribution.

We evaluate OOD detection performance on our OOD detection testing dataset, containing 1,050 OOD prompts and 954 in-distribution prompts (See Section \ref{sec:dataset_prep}).
% We use True Positive Rate (TPR) to measure detection accuracy, defined as the proportion of OOD prompts correctly identified as OOD:
% $\text{TPR} = \frac{\text{\# correctly detected OOD prompts}}{\text{\# total OOD prompts}}$.
% We use False Positive Rate (FPR) to quantify the frequency of misclassifying in-distribution prompts as OOD, defined as $\text{FPR} = \frac{\text{\# in-distribution prompts misclassified as OOD}}{\text{\# total in-distribution prompts}}$. 
We measure performance using precision, recall, and F1-Score, where correctly detected OOD prompts are treated as true positives. Precision measures the proportion of correctly identified OOD prompts among all prompts flagged as OOD, recall measures the proportion of actual OOD prompts that were correctly detected, and F1-Score provides a balanced measure combining both precision and recall.
% \mike{Ray to double check whether the definition of these metrics are correct}

%We then run inference the OOD Detection Testing Dataset in Section \ref{sec:dataset_prep}, consisting of 1,039 OOD samples and 954 in-distribution samples.
%We use the True Positive Rate (TPR) to quantify the defensive capability of guardrails, defined as the percentage of the number of the jailbreak prompts that can be the successfully detected as OOD by a runtime guardrail $\#\mathrm{Prompts}_\mathrm{success}$ and the total number of the jailbreak prompts $\#\mathrm{TotalPrompts}$:
%\begin{equation}
%\mathrm{TPR} = \frac{\#\mathrm{Prompts}_\mathrm{success}}{\#\mathrm{TotalPrompts}}
%\label{DSR_formula}
%\end{equation}
%In this case, we define successfully defended by if they are identified as OOD at each threshold value, for each metric.
%The False Positive Rate (FPR) is calculated by the percentage of prompts from the validation dataset that are identified as OOD at each threshold.

\subsubsection*{\underline{\textbf{Results}}}
% \input{tables/RQ1_OOD_detection}

\begin{figure}
    \centering
    \includegraphics[width=1\linewidth]{figs/RQ1_F1.pdf}
    \vspace{-0.3cm}
    \caption{(RQ1 Results) OOD detection performance of each OOD method evaluated in RQ1.}
    \label{fig:RQ1_OOD}
\end{figure}

% \begin{figure*}
%     \centering
%     \includegraphics[width=1\linewidth]{figs/RQ1_ood_detection_combined_vs_percentile.png}
%     \caption{Combined OOD detection and False Positive Rate ofRQ1}
%     \label{fig:RQ1_FPR}
% \end{figure*}

Fig. \ref{fig:RQ1_OOD} presents the Precision, Recall, and F1-Score of the four OOD detection methods, evaluated under three quantile thresholds (90th, 95th, 99th).

%\textbf{The best balanced performance was Mahalanobis Distance with 95.48\% TPR and 3.50\% FPR, at the 99th quantile threshold.}
\textbf{The Mahalanobis Distances used in our framework achieves the highest F1-Score of 96.1\% when the OOD detection threshold $\tau\_{\text{OOD}}$ is set to the 99th quantile threshold.}
Across all methods, increasing the detection threshold $\tau\_{\text{OOD}}$ generally improves precision while reducing recall. This is because a higher threshold makes the detector more conservative, flagging fewer inputs as OOD, which reduces false positives (improving precision) but may miss some actual OOD samples (reducing recall).
%In particular, the Mahalanobis distance achieves the highest recall of 98.7\% at the 90th percentile threshold, but with only 84.4\% precision, demonstrating how lower thresholds favor recall at the expense of precision.
Notably, the Likelihood Ratios and Energy Scores present less balanced performance, with their precision-recall trade-offs resulting in lower F1 scores across all thresholds.
% This represents an 8\% improvement over the best-performing baseline, the energy score, at the same threshold.
On the other hand, the balanced F1-Score of 96.1\% confirms that Mahalanobis Distances provides the optimal trade-off between precision and recall.
These results confirm the feasibility of leveraging OOD detection to enable guardrails to recognize unseen inputs continually. Among all methods, the Mahalanobis Distances proves to be the most effective in identifying OOD samples while maintaining a low false positive rate.
%The TPR of all OOD detection methods except Ensemble Uncertainty achieved 84\%+ across 90\%, 95\% and 99\% quantile threshold values.
%For instance, the Mahalanobis distance maintains high TPR (96\% - 99\%), while its FPR reduced from 20\% to 3.5\% at 99th quantile.
%Similarly, Energy Scores achieved a strong balanced performance of around 88\% TPR at all quantiles, and the FPR dropped from 11.1\% to 1.2\%.
%In contrast, Ensemble Uncertainties showed a significant drop in TPR, from 88\% at 90th quantile, to 8.1\% at 95th quantile, then only 1.3\% at 99th quantile. 
%Thess results highlight that while these methods are able to identified unknown jailbreak attacks at a high ratio, threshold selection directly impacts FPR and sometimes can heavily impact TPR too.
%However, the Energy Score and Mahalanobis Distance metrics demonstrated superior stability and effectiveness under most quantile selections.
%This finding confirms our hypothesis that OOD detections methods are capable of detecting unusual jailbreak prompts that deviate from expected prompts in the training dataset.

\subsection{RQ2: \rqtwo}
\subsubsection*{\underline{\textbf{Approach}}}
\begin{figure*}[h]
    \centering
    \includegraphics[width=\textwidth]{figs/RQ2_DSR_Comparison.pdf}
    \caption{(RQ2 Results) The performance comparison of our \ourapp~and LlamaGuard on OOD prompts when continuously updated through detected OOD prompts.}
    \label{fig:RQ2_DSR}
\end{figure*}
% talk about for each new attack we use the base model again
To address this RQ, we compare our \ourapp~approach with LlamaGuard~\cite{inan2023llama}, a state-of-the-art open-source runtime guardrail for LLM-powered systems. Although LlamaGuard shows strong performance on in-distribution data, our analysis reveals its limitations in handling unsafe prompts that differ significantly from its training distribution. We use the LlamaGuard-3-1B and LlamaGuard-3-8B models from Hugging Face as baselines.

To simulate continual learning under OOD detection, we construct a Continual Learning Testing Dataset of 992 OOD jailbreak prompts, all correctly identified by our method in RQ1 (see Section~\ref{sec:dataset_prep}). These prompts span ten types of jailbreak attacks, each exhibiting distinct patterns that bypass guardrails trained on natural unsafe prompts, making them valid OOD scenarios.

In real-world settings, jailbreaks tend to emerge chronologically as adversaries innovate new attack patterns. To reflect this, we organize the dataset into sequential attack waves, where each wave represents a set of related jailbreak techniques. Instead of training on all attacks simultaneously, models are updated incrementally.

For each wave, we hold out 50 samples for testing and use the remaining 50 OOD samples for continual learning. Prompts are fed to the model one at a time. At each time step, we first run inference to log the model’s prediction, then update the model using the same sample. This ensures no data leakage between training and evaluation. The process continues iteratively until all non-hold-out samples are used.

We evaluate two guardrails: our trained \ourapp~from RQ1 and the LlamaGuard baselines.
For each attack wave, all models are initialized from the same pre-trained checkpoint to ensure a fair and consistent comparison.
All models are updated using the same LoRA-based continual learning procedure after running inference.
We measure each guardrail’s effectiveness using the Defense Success Rate (DSR), defined as the number of jailbreak prompts correctly predicted as unsafe, divided by the total number of jailbreak prompts.
This setup allows us to assess how rapidly and effectively each guardrail adapts to unseen OOD jailbreaks using limited adaptation data.


%To address this RQ, we investigate the effectiveness of LoRA-based continual fine-tuning for updating guardrails using jailbreak prompts identified as OOD from RQ1. We apply LoRA-based fine-tuning to both the pre-trained~\ourapp from RQ1, as well as LlamaGuard 3, a state-of-the-art guardrail, following the procedures in Section \ref{sec:exp_design}. 

%In each jailbreak attack, we hold out 50 samples as the testing set, and update the guardrails using the 50 out of the remaining samples that are identified as OOD. To avoid any overlap between models, each attack will use the base pre-trained~\ourapp in RQ1, and base LlamaGuard model, at the start of the finetuning process. Each LoRA-based adaptation iteration is conducted with one jailbreak prompt only, and evaluated on the 50 hold-out testing set, allowing us to assess how~\ourapp learns with limited data. This process continues for 50 iterations or until the training dataset is exhausted. This setup enables us to observe how quickly and effectively each guardrail adapts to new OOD attacks with limited data using LoRA-based continual learning.

%In an actual production environment, jailbreak attacks emerge chronologically as adversaries continuously develop new techniques to bypass model safeguards. 
%To simulate this real-world scenario, we structure our experiment to process jailbreak attacks in temporal batches. We organise our OOD jailbreak examples into sequential 'wave' of attack, with each wave representing a collection of related jailbreak techniques that might emerge during a specific time period. 
%Rather than training on all available data at once, we incrementally expose our models to these attack waves, allowing us to evaluate how each updating technique adapts to newly emerging threats.

%For each attack wave, we apply our four updating methods independently and measure their ability to incorporate new defence capabilities without compromising existing ones. 
%This approach closely mirrors production deployment scenarios where guardrail models must be continuously updated against evolving attack strategies while maintaining consistent performance across all previously encountered threats.

% For SFT on~\ourapp, we first conducted detailed parameter sweep to identify the best hyper parameter to use on a separate jailbreak dataset. To name a few: 5e-5 learning rate, lambda value of 0.005, max gradient norm of 1. These settings were then fixed for all subsequent SFT experiments to ensure consistency and comparability.

% For LoRA-based finetuning, we adopted the commonly used LoRA configuration, setting the rank to 8, the alpha parameter to 32, and a dropout rate of 0.1. These values are consistent with best practices reported in recent literature and ensure a fair comparison with SFT and full-parameter finetuning \cite{hu2022lora}.


\subsubsection*{\underline{\textbf{Results}}}
Fig. \ref{fig:RQ2_DSR} presents the time-wise DSR across each attack wave, showing how \ourapp, LlamaGuard-1B, and LlamaGuard-8B defend against OOD prompts over time.
We also present three additional baselines where the continual learning (CL) is not applied.
Each subplot represents an attack wave. The x-axis indicates the number of consecutive update steps, while the y-axis shows the DSR on the 50 held-out unsafe samples for each attack.

\textbf{Our~\ourapp~+ Continual Learning (CL), consistently shows the fastest adaptation across most attack waves, achieving optimal DSR within 2 to 38 update steps, with a median of 2 update steps.
In comparison, LlamaGuard requires between 4 to 44 update steps to achieve optimal DSR across the same attack waves, with a median of 4 steps. Guardrails that do not incorporate our continual learning (CL) framework maintain the same level of performance across time steps for each attack wave, as they are not updated with detected OOD prompts.}
\textbf{These results confirm that our CL framework is crucial for enabling adaptive defenses, allowing the guardrails to continually improve and remain effective against evolving and previously unseen jailbreak attacks.}

We also found that different attack patterns exhibited varying levels of complexity for the models to adapt to, significantly impacting their adaptation speed.
In particular, template based attacks such as AIM and SmartGPT converged to optimal DSR in only 2 updating steps for \ourapp~+ CL, and around 4-6 steps for LlamaGuard 8B + CL. 
On the other hand, attacks like Base64 and Zulu presented more significant challenges. Both ~\ourapp + CL and LlamaGuard + CL presented notably slower adaptation speed and often plateaued at optimal DSR values lower than those of less sophisticated attacks.
This indicates that these more complex jailbreak attacks require more extensive adaptation.
%In nearly every scenario, AdaptiveGuard + LoRA reaches 100\% DSR within the first 10--20 iterations,
% \textbf{These results demonstrate the effectiveness of our \ourapp~at learning to defend against new jailbreak attacks with a small number of OOD samples.}
%This performance is competitive with, and in many cases matches or exceeds, the LlamaGuard 1B + LoRA and LlamaGuard 8B + LoRA baselines.


% As a baseline, \ourapp~without CL has a median DSR across 10 attacks of 81\%. With CL, its performance reaches a median DSR of 100\% (a 19\% improvement).
% LlamaGuard-1B performs similarly, with a median DSR of 82\% without CL, improving to 100\% with CL. The larger LlamaGuard-8B model starts with a much lower median DSR of 57\%, but like the others, reaches a perfect median DSR of 100\% when CL is applied.

% \ray{discuss three baseline separately, medium performance, and then after adopting CL -> how much \% improvement}
% All three baselines without CL exhibit DSRs ranging from 0\% to 100\%, with a median performance of 74\%. 
% In contrast, our \ourapp~approach, when updated with detected OOD prompts, achieves substantially higher DSRs—from 90\% to 100\%, with a median of 100\%.

% \ray{remove 100\% change into optimum}
% \ray{how much faster is our compare to llamaguard}
% \ray{adaptation speed based on different attacks, what attacks are difficult}
% \ray{intro, RQ2 motivation, RQ2 result}

%In addition, the guardrails that are not updated with the detected OOD prompts, shown in the dashed lines, display varying results across different attack types.
%AdaptiveGuard with no updates had 100\% DSR from the start against attacks such as Caesar Cipher, Combination and Self Cipher, while performed poorly against Code Chameleon attack with 0\% DSR. The LlamaGuard 8B model with no updates showed DSR ranging from 17\% to 82\%, while the 1B model had similar results but achieved 100\% DSR in Caesar Cipher and Code Chameleon attacks.

%There is some variation in learning speed and final DSR across different attack types. For example, attacks such as AIM, DAN, SmartGPT and DeepInception were quickly adapted towards 100\% DSR, while a few attacks (e.g., Base84, Zulu) required more iterations to reach similar DSR. 

\subsection{RQ3: \rqthree}

\begin{figure*}[htbp]
    \centering
    \includegraphics[width=\textwidth]{figs/RQ3_F1_Comparison.pdf}
    \caption{(RQ3 Results) The performance comparison of our \ourapp~and LlamaGuard on in-distribution prompts when continuously updated through detected OOD prompts.}
    \label{fig:RQ3_F1}
\end{figure*}

\subsubsection*{\underline{\textbf{Approach}}}
To answer this RQ, we reuse the experimental setup from RQ2 for continual learning (CL), with one key difference: after each continual update step, we evaluate on in-distribution (ID) data instead of out-of-distribution (OOD). This allows us to assess knowledge retention of ID prompts after updating with each OOD prompt. Specifically, we use the Aegis validation set—part of the dataset from RQ1—containing 314 safe and 640 unsafe natural language prompts. We measure performance using the F1-score to account for both safe and unsafe classes.

%To answer this RQ, we reuse the experimental setup in RQ2 for continual learning (CL).
%We evaluate the guardrails on the Aegis validation dataset (in-distribution prompts) used in RQ1 after each continuous update step to test the knowledge retention of each guardrail.
%This dataset includes natural safe and unsafe prompts that are considered in-distribution. 

%represents the original task the guardrails were trained for, and high performance on this set indicates successful knowledge retention.

%Specifically, for each attack type, the guardrails are updated iteratively using OOD jailbreak samples, as described in RQ2. At each iteration, we record the performance on the original safe/unsafe validation set. This allows us to track how well each updating method preserves the guardrails' ability to correctly classify previously learnt types of prompts, even as it learns to defend against new attacks. The details of the dataset can be found in Section \ref{sec:dataset_prep}.

\subsubsection*{\underline{\textbf{Results}}}
%Figure~\ref{fig:RQ3_F1} presents the change in validation F1 score for each attack type and updating method over 50 adaptation iterations.
Fig.~\ref{fig:RQ3_F1} presents the time-wise F1-Score across each attack wave, showing how \ourapp, LlamaGuard-1B, and LlamaGuard-8B defend against in-distribution prompts over time when continuously updated with OOD prompts.
Similar to RQ2, we present three additional baselines where the continual learning (CL) is not applied.
Each subplot represents an attack wave. The x-axis indicates the number of consecutive update steps, while the y-axis shows the F1-Score on the 954 in-distribution prompts from the Aegis validation set.

%50 held-out unsafe samples for each attack.
%Each subplot corresponds to a specific attack, and the y-axis represents the F1 score on the original validation set, reflecting the model’s ability to retain its in-distribution classification performance as it adapts to new OOD jailbreak prompts.
% \textbf{The results reveal that our proposed~\ourapp~approach achieves strong adaptation to new attacks while maintaining high retention of prior knowledge.} 
\textbf{\ourapp~+CL achieves the highest median F1-Score of 85\% on in-distribution prompts across all attacks after the final update with OOD prompts, with the F1-Score ranging from 85\% to 86\%.}
In comparison, the best-performing baseline, LlamaGuard-8B, achieves a median F1-Score of 80\% on in-distribution prompts, with a minimum of 77\% and a maximum of 81\%.

Across all attack types, \ourapp~+ CL consistently maintains the highest in-distribution F1-Score throughout the CL process, with only slight variation (ranging from –0.4\% to +0.4\%).
This indicates that \ourapp~+CL is able to learn to defend against unseen jailbreak attacks without sacrificing its learned in-distribution knowledge.
\textbf{In summary, \ourapp~+ CL maintains the highest performance on in-distribution prompts compared to all baselines, demonstrating minimal forgetting even after continuous updates with detected OOD prompts.}
% This suggests that LoRA-based method on GPT2Guard are less prone to catastrophic forgetting, allowing the model to learn new attacks while preserving its original capabilities.

\section{Discussion}
\label{sec:discussion}
In the previous section, we evaluated the effectiveness of our proposed \ourapp~approach for OOD detection and continual learning against other baselines.
While our results demonstrate clear advancements, it remains unclear (1) whether our approach can still perform well in a setting where a single continual model is built across all attacks, rather than reinitialising the model for each attack; (2) what the trade-offs are between different continual learning (CL) update methods; and (3) what is the computational efficiency of our approach compared to other baselines. Thus, we evaluate a single continual model across all attacks, compare LoRA with full SFT as update strategies, and analyze the computational efficiency of each approach.
%While our results demonstrate clear advancements, there exist several practical considerations to further analyse.
%In this section, we examine three key aspects: (1) knowledge retention of previously learnt attacks in sequential continual learning scenarios, (2) the trade-offs between different guardrail updating methods (LoRA vs. SFT), and (3) computational efficiency compared to baseline approaches. These discussions provide insights into the practical deployment implications and limitations of our approach.

\begin{table}[h]
\centering
\caption{(Discussion) DSR of \ourapp~under Sequential Continual Learning with and without Early Stopping across 10 OOD Attack Types.}
\label{tab:lora_dsr_results}
\resizebox{\linewidth}{!}{%
\begin{tabular}{lcc}
\toprule
\textbf{Attack Type} & \textbf{DSR (No Early Stopping)} & \textbf{DSR (Early Stopping)} \\
\midrule
AIM Attack & 1.00 & 1.00 \\
Base64 Attack & 1.00 & 0.86 \\
Caesar Attack & 1.00 & 1.00 \\
Code Chameleon Attack & 1.00 & 1.00 \\
Combination Attack & 1.00 & 1.00 \\
DAN Attack & 1.00 & 1.00 \\
DeepInception Attack & 1.00 & 1.00 \\
Self Cipher Attack & 1.00 & 1.00 \\
SmartGPT Attack & 1.00 & 1.00 \\
Zulu Attack & 1.00 & 0.92 \\
\midrule
\textbf{Average} & \textbf{1.00} & \textbf{0.98} \\
\bottomrule
\end{tabular}
}
\end{table}

\subsection{Continual Performance Without Model Reinitialization}
%\ray{say why we do the setup in RQ2 and RQ3 (fair comparison between different methods), but in the industry's context, we only have 1 continuous model, so in the context, what would it look like }
%In RQ2 and RQ3, we initialised a separate copy of the base-trained~\ourapp~from RQ1 for each attack. This approach ensured a fair comparison by isolating the effects of individual attacks on model adaptation performance.
%In industry settings, guardrails are typically deployed as a single, continuously updated model that incrementally incorporates new external knowledge over time.
%In that context, while RQ3 focuses on the evaluation of~\ourapp with whether it has retained previous knowledge on classifying the validation dataset, it does not explicitly evaluate whether learning new attacks does not compromise knowledge retention of previously learnt attacks.
In a typical industry setup, a guardrail is deployed as an additional layer before user inputs reach the underlying LLM. This guardrail often operates as a single continually updated model that incrementally incorporates new external knowledge—such as OOD prompts—to adapt over time. To reflect this real-world scenario, we construct a single continual version of \ourapp, which is updated sequentially across all 10 attacks without reinitialization between them.

We apply continual learning by updating the same \ourapp~model using OOD prompts from each attack in the same experimental dataset used in RQ2. After each update, we evaluate the model's DSR on the test set corresponding to the respective attack. We consider two update settings: in the early stop setting, the continual learning process halts once the model achieves 95\% DSR on the current attack; in the no early stop setting, training continues through all OOD prompts regardless of intermediate performance.

Table~\ref{tab:lora_dsr_results} presents the DSR results for our sequential continual learning setup. The results show that \ourapp~achieves a 100\% DSR across all 10 attacks when trained without early stopping. Even with early stopping enabled, the model maintains an average DSR of 98\%.
Notably, the Obfuscation-Based attacks—Base64 (86\%) and Zulu (92\%)—prove the most challenging to defend under early stopping. This suggests that these attacks require more adaptation steps for the model to fully internalize their patterns in a guardrail.
Overall, these results indicate that our continual learning framework not only adapts effectively to new attack types but also retains knowledge of previously encountered ones. This makes it well-suited for real-world deployment, where a single guardrail must continuously evolve to defend against emerging threats without forgetting past vulnerabilities.
%\ray{Therefore, we conducted additional experiments to evaluate this by continuously updating a single~\ourapp~model with all 10 attacks sequentially, and then evaluating its performance on the test set of each individual attack. Table~\ref{tab:lora_dsr_results} shows the DSR of this sequential continual learning approach compared to the isolated training approach from RQ2. The results demonstrate that our LoRA-based~\ourapp~achieves perfect defense success rates (DSR = 1.00) across all 10 attacks when trained without early stopping, exceeding the DSR achieved in RQ2. Even with early stopping enabled, the model achieves an average DSR of 0.98, with only minor performance degradation on Base64 (0.86) and Zulu (0.92) attacks. This indicates that our LoRA-based continual learning approach effectively preserves knowledge of previously learnt attacks while adapting to new ones, making it suitable for real-world deployment scenarios where guardrails must continuously evolve to defend against emerging threats.}

\begin{table}[htbp]
\centering
\caption{(Discussion) Comparison of SFT vs LoRA Performance on Jailbreak Defense}
\label{tab:sft_lora_comparison}
\begin{tabular}{lcccc}
\toprule
\multirow{2}{*}{\textbf{Statistic}} & \multicolumn{2}{c}{\textbf{Full Fine-Tuning}} & \multicolumn{2}{c}{\textbf{LoRA}} \\
\cmidrule(lr){2-3} \cmidrule(lr){4-5}
& \textbf{DSR} & \textbf{F1-Score} & \textbf{DSR} & \textbf{F1-Score} \\
\midrule
Min & 1.00 & 0.800 & 0.90 & 0.848 \\
Median & 1.00 & 0.800 & 1.00 & 0.852 \\
Max & 1.00 & 0.801 & 1.00 & 0.856 \\
\midrule
\textbf{Average} & \textbf{1.00} & \textbf{0.800} & \textbf{0.99} & \textbf{0.852} \\
\bottomrule

\end{tabular}
\end{table}


\subsection{Trade-offs Between LoRA and Full Fine-Tuning for CL}
In our continual learning (CL) framework, we adopted LoRA-based adaptation for updating the guardrails.
LoRA is known for its computational efficiency, as it updates only a small subset of parameters while keeping the rest of the pre-trained model frozen. This selective adaptation can lead to better knowledge retention—reflected as stronger in-distribution performance in our context—and has been shown to mitigate catastrophic forgetting compared to full-parameter fine-tuning~\cite{hu2022lora,he2022towards}. While our results confirm LoRA's effectiveness, it remains unclear whether it consistently offers better knowledge retention than full fine-tuning in the context.
%LoRA (Low-Rank Adaptation) offers several practical advantages: it is computationally efficient, requires updating only a small subset of parameters, and has been shown to mitigate catastrophic forgetting compared to full-parameter fine-tuning in many settings~\cite{hu2022lora,he2022towards}. Our results demonstrate that LoRA enables rapid adaptation to new jailbreak attacks while maintaining strong knowledge retention on previously learnt tasks.

%However, it is important to acknowledge that different fine-tuning strategies may yield different trade-offs between adaptation speed, model performance, and knowledge retention. For example, Supervised Fine-Tuning (SFT) \cite{howard2018universal}, which updates all model parameters, may allow for better learning of new threats but at the risk of increased forgetting of prior knowledge, as observed in prior continual learning literature~\cite{parisi2019continual}. 

To investigate this, we compare our LoRA-based approach~\cite{hu2022lora} with full fine-tuning~\cite{howard2018universal} using the same experimental setup as in RQ2 and RQ3. Since the evaluation spans ten different attacks, we report the minimum, median, and maximum values for both DSR (representing performance for OOD prompts) and F1-Score (performance for in-distribution prompts).
%
Table~\ref{tab:sft_lora_comparison} presents the performance comparison between LoRA and full fine-tuning across all 10 attacks. Our LoRA-based approach achieves a median DSR of 100\% and a median F1-Score of 85\% on in-distribution prompts, indicating strong knowledge retention. In contrast, full fine-tuning also achieves a median DSR of 100\% but yields a lower median F1-Score of 80\%. These results are consistent with prior findings in the continual learning literature and demonstrate that our LoRA-based method more effectively preserves previously learned knowledge while adapting to new threats.


% Future work should systematically compare a broader range of adaptation techniques for continual learning framework. Such studies would provide deeper insights into the optimal strategies for guardrail updating in real-world LLM-powered systems, where both rapid adaptation and robust knowledge retention are essential for long-term safety and reliability.

\begin{table}[h]
\centering
\vspace*{-0.5em}
\caption{(Discussion) Training \& inference time per sample, and memory usage during training}
\label{tab:computation_requirements}
\setlength{\tabcolsep}{4pt}
\begin{tabular}{lccc}
\toprule
\textbf{Metric} & \textbf{AdaptiveGuard} & \textbf{LlamaGuard 1B} & \textbf{LlamaGuard 8B} \\
 & \textbf{LoRA} & \textbf{LoRA} & \textbf{LoRA} \\
\midrule
Training Time & 0.60s & 1.06s & 2.04s \\
Inference Time & 0.01s & 0.25s & 1.10s \\
Memory Usage & 1.3 GB & 4.0 GB & 27.1 GB \\
\bottomrule
\vspace*{-1.5em}
\end{tabular}
\end{table} 

\subsection{\hbox{Computational Efficiency: \ourapp~and Baselines}}
In the context of practical deployment of our \ourapp, computational resources are a key consideration factor.
Therefore, we analyse (1)~the training time per iteration (update step), (2)~inference time per testing sample, and (3)~training and testing memory usage.
We conduct the analysis using the datasets from RQ2 and RQ3, comprising 992 jailbreak attacks, 314 safe, and 640 unsafe NL prompts.
We compare our approach (\ourapp+LoRA) with the baselines LlamaGuard-1B+LoRA and LlamaGuard-8B+LoRA.
%We measure time in milliseconds and memory in gigabytes.

%When compared to the full SFT of ~\ourapp, the~\ourapp+LoRA CL approach reduces training time by 20\% (0.60 vs 0.75 seconds per iteration) while requiring 57\% less memory (1.3 GB vs 3.0 GB). 
%Most notably, LoRA updates only 17.08\% of the model parameters (21.3M) compared to the 100\% parameter updates (124M) required for full fine-tuning, making our approach substantially more parameter-efficient.
Table \ref{tab:computation_requirements} summarises the computational efficiency of our approach compared to baseline methods during both training and inference.
Compared to the LlamaGuard baselines, our method demonstrates substantially greater efficiency.
Specifically, compared to LlamaGuard-1B and LlamaGuard-8B, \ourapp~achieves 43\% and 71\% faster training times during continual learning (CL), delivers 25× and 110× faster inference, and reduces memory usage by 67\% and 95\%, respectively.
These substantial gains stem from the compact architecture of our base model, GPT-2 (approximately 137M parameters), in contrast to the considerably larger LlamaGuard models with 1B and 8B parameters.
These results highlight the practicality of our approach for resource-constrained settings.

%\input{sections/related_work}

\section{Threats to Validity}
\label{sec:threats}
\textbf{Threats to construct validity} relate to the selection of jailbreak attacks and OOD thresholds.
We selected 10 different jailbreak attacks guided by their prevalence to guardrails and ability to represent a wide spectrum of jailbreak techniques \cite{dong2024safeguarding, yi2024jailbreak, wei2024jailbroken}. 
While additional attack types could be included in future evaluations, this would not fundamentally alter the key conclusions presented in our research questions: that complex jailbreak attacks can be effectively detected by our proposed OOD detection measures (RQ1), and that detected attacks can be efficiently learned by \ourapp through continual learning with minimal catastrophic forgetting (RQ2 and RQ3).
%
For the OOD threshold, we adopted commonly used quantile-based thresholds (90\%, 95\%, and 99\%) based on the distribution of OOD scores on the training set~\cite{liu2020energy,lee2018simple}. This approach allows us to select the most appropriate threshold values based on our experiments, ensuring that only the most anomalous prompts are flagged for guardrail adaptation. However, in practice, threshold selection is based on various factors depending on the application's tolerance for risk and the distributional properties of the data. 

\textbf{Threats to internal validity} relate to the potential influence of hyperparameter settings during the fine-tuning of our \ourapp~and the selection of OOD detection thresholds. Variations in energy threshold values, LoRA configurations, or learning rates, compared to those specified in Section \ref{sec:exp_design}, could impact experiment's outcomes. 
To address these threats, we open-source our replication package and provide detailed documentation of all hyperparameter settings to ensure the experiment is reproducible by future researchers. 

\textbf{Threats to external validity} concern the generalizability of our results. Our experiment findings are supported by the dataset, jailbreak methods, and model architecture employed during the study. Our training dataset contains 12,233 unsafe prompts and 6,905 safe prompts. Our  validation dataset contains 314 safe and 640 unsafe natural language prompts to assess false positive rates and catastrophic forgetting of~\ourapp. We also applied the 10 jailbreak attacks to 105 separate unsafe prompts, resulting in 1,050 jailbreak prompts across 10 categories for RQ1. Of the 1,050 jailbreak prompts, 992 are identified as OOD and used to study the continual learning adaptation in RQ2 and RQ3.
Our evaluation focuses on GPT-2 as the base model, which represents a widely-studied language model in AI safety research, but may not capture the behaviors of larger or more recent language models.
While \ourapp~is fine-tuned specifically to address the jailbreak attacks and model configurations discussed in this paper, other prompt datasets, jailbreak methods, and language model architectures can be explored in future work.


% \subsection{OOD Threshold Selection}
%  In our experiments, we adopted commonly used quantile-based thresholds (90\%, 95\%, and 99\%) based on the distribution of OOD scores on the training set~\cite{hendrycks2017baseline,liu2020energy,lee2018simple}. This approach allows us to select the most appropriate threshold values based on our experiments, ensuring that only the most anomalous prompts are flagged for guardrail adaptation. However, in practice, threshold selection are based on various factors depending on the application's tolerance for risk and the distributional properties of the data. Future works should focus on the different threshold selection methods to achieve the optimal performance on different metrics \cite{lipton2014optimal}. 

\section{Conclusion}
In this paper, we first show that state-of-the-art guardrails like LlamaGuard face a critical limitation, with their Defense Success Rate (DSR) dropping to as low as 12\% against jailbreak attacks not seen during training.
To address this, we present \ourapp, an OOD-aware continual learning framework that detects and adapts to previously unseen jailbreak patterns by treating them as out-of-distribution (OOD) inputs. This builds on the observation that guardrails like LlamaGuard are trained on NL unsafe prompts, while jailbreaks often use obfuscated inputs that fall outside their training distribution.
To evaluate \ourapp, we compile a dataset using ten state-of-the-art jailbreak methods: AIM, DAN, Self Cipher, Deep Inception, SmartGPT, and Code Chameleon.
Through our evaluation, we found that \ourapp~achieves a 96\% true positive rate in OOD detection and reaches 100\% DSR within a median of two update steps—twice as fast as LlamaGuard under the same continual learning setup.
Moreover, \ourapp~retains 85\% F1-score on in-distribution data after adaptation, outperforming LlamaGuard’s 80\%. These results demonstrate that \ourapp~could offer an effective post-deployment solution for adaptive jailbreak defense in dynamic production.
%In this paper, we demonstrate that state-of-the-art guardrails such as LlamaGuard face critical vulnerabilities with Defence Success Rate (DSR) dropping to as low as 12\% under unseen jailbreak attacks. To address this, we introduce~\ourapp, an OOD-aware continual learning framework that detects and adapts to unseen jailbreak patterns. \ourapp~achieves 96\% True Positive Rate for OOD detection and reaches 100\% DSR within a median of 2 update steps, compared to LlamaGuard's 4 steps, while maintaining 85\% F1-score on in-distribution data with minimal catastrophic forgetting. Our results establish~\ourapp~as an effective solution for deploying adaptive guardrails in production environments where security threats continuously evolve.

%\appendix
%\section{Appendix}
%...

% Bibliography
\balance
\bibliographystyle{IEEEtran}
\bibliography{reference}

\end{document}
